
This section defines the physical system, the baseline control architecture, the cable-snap fault event, and the formal problem.

% ============================================================
\subsection{System Model}
% ============================================================

\subsubsection{Configuration Space}

The system comprises $N$ quadrotors cooperatively transporting a rigid payload via flexible cables. The state evolves on
\begin{equation}\label{eq:config_manifold}
  \mathcal{Q} = \SEthree \times \bigl(\SEthree \times \Sph^2\bigr)^N, \quad \dim(\mathcal{Q}) = 6 + 8N,
\end{equation}
where $\SEthree$ captures the payload's six-DOF pose, each quadrotor contributes a further $\SEthree$, and each cable direction lives on $\Sph^2$. The default configuration uses $N = 3$ ($\dim(\mathcal{Q}) = 30$), with scaling experiments extending to $N \in \{5, 7\}$.

\subsubsection{Quadrotor Dynamics}

Each quadrotor $i \in \{1,\dots,N\}$ has mass $m_Q = 1.5$~kg and obeys (with $\bm{e}_3 = [0,0,1]^\top$ pointing upward and $g = 9.81$~m/s$^2$):
\begin{align}
  m_Q \ddot{\bm{p}}_i &= -m_Q g\,\bm{e}_3 + f_i R_i \bm{e}_3 + \bm{F}_i^{\mathrm{cable}} + \bm{F}_i^{\mathrm{wind}}, \label{eq:quad_trans} \\
  J \dot{\bm{\Omega}}_i &= -\bm{\Omega}_i \times (J\bm{\Omega}_i) + \bm{\tau}_i + \bm{\tau}_i^{\mathrm{ext}}, \label{eq:quad_rot} \\
  \dot{R}_i &= R_i \hatmap{\bm{\Omega}_i}, \label{eq:quad_kin}
\end{align}
where $f_i \in [0, f_{\max}]$ is the scalar thrust ($f_{\max} = 150$~N), $R_i \in \SOthree$ is the body-to-world rotation, $J \in \R^{3 \times 3}$ is the inertia tensor (from a $0.30 \times 0.30 \times 0.10$~m solid box), $\bm{\Omega}_i \in \R^3$ is the body angular velocity, $\bm{\tau}_i \in \R^3$ is the commanded torque ($\abs{\tau_k} \leq 10$~N$\cdot$m per axis), $\bm{F}_i^{\mathrm{cable}}$ is the cable force, and $\bm{F}_i^{\mathrm{wind}}$ is the aerodynamic disturbance.

\subsubsection{Payload Dynamics}

The payload is a rigid sphere of mass $m_L = 3.0$~kg and radius $r_L = 0.15$~m:
\begin{equation}\label{eq:payload_dynamics}
  m_L \ddot{\bm{p}}_L = -m_L g\,\bm{e}_3 + \textstyle\sum_{i=1}^{N} T_i \bm{q}_i + \bm{F}^{\mathrm{contact}} + \bm{F}_L^{\mathrm{wind}},
\end{equation}
where $T_i \geq 0$ is the scalar cable tension, $\bm{q}_i \in \Sph^2$ is the unit direction from payload to quadrotor~$i$, and $\bm{F}^{\mathrm{contact}}$ accounts for ground interaction (Coulomb friction, $\mu_s = 0.5$, $\mu_k = 0.3$). Cable attachment points are at radius $0.105$~m on the payload surface. The payload mass $m_L$ is \emph{unknown to every agent}.
The force balance in \eqref{eq:payload_dynamics} is the plant-level relation that later explains how cable-load redistribution appears immediately at the payload during a snap event.

\subsubsection{Bead-Chain Cable Model}

Each cable is modeled as a chain of $n_b = 8$ point-mass beads ($m_{\mathrm{rope}} = 0.2$~kg) creating $n_b + 1 = 9$ spring-damper segments with tension
\begin{equation}\label{eq:cable_tension}
  \mathcal{T}_j = \begin{cases}
    0, & \text{if } \epsilon_j \leq 0 \;\text{(slack)}, \\
    k_s \epsilon_j + c_s \max(0, \dot{\epsilon}_j), & \text{if } \epsilon_j > 0 \;\text{(taut)},
  \end{cases}
\end{equation}
where $\epsilon_j = \norm{\bm{b}_{j-1} - \bm{b}_j} - L_0$ is the segment stretch and $L_0 = L_{\mathrm{rest}}/(n_b + 1)$. Stiffness is derived from maximum allowable strain $\varepsilon_{\max} = 0.15$:
\begin{equation}\label{eq:cable_stiffness}
  k_s = \frac{m_L g}{N \cdot L_{\mathrm{rest}} \cdot \varepsilon_{\max}} \cdot (n_b + 1).
\end{equation}
The piecewise law in \eqref{eq:cable_tension} is the key cable nonlinearity for the reported fault transients, because it captures the slack-to-taut switching that a rigid-cable model would miss.
Damping is $c_s = c_{\mathrm{ref}} \sqrt{k_s / k_{\mathrm{ref}}}$ with $c_{\mathrm{ref}} = 15$~N$\cdot$s/m, $k_{\mathrm{ref}} = 300$~N/m. Cable lengths are sampled from per-quadrotor Gaussians: $L_i \sim \mathcal{N}(\mu_i, s_i^2)$ with means $\mu \in \{1.0, 1.1, 0.95\}$~m and standard deviations $s \in \{0.05, 0.08, 0.06\}$~m, producing up to $19\%$ asymmetry.

The tension-only constraint means cables cannot push. The bead-chain discretization captures lateral vibration, wave propagation ($v_{\mathrm{wave}} \approx 7$~m/s), distributed inertia, and slack-to-taut transitions that rigid-cable assumptions miss. The plant is simulated in Drake~\cite{drake2024} at $\Delta t_{\mathrm{sim}} = 0.2$~ms ($5$~kHz).

\subsubsection{Wind Disturbance}

Aerodynamic disturbance follows the MIL-F-8785C Dryden turbulence model ($\sigma_u = \sigma_v = 1.0$~m/s, $\sigma_w = 0.5$~m/s; $L_u = L_v = 200$~m, $L_w = 50$~m; mean wind $[2.0, 0, 0]$~m/s). Spatial correlation decays as $\rho(\bm{p}_i, \bm{p}_j) = \exp(-\norm{\bm{p}_i - \bm{p}_j}/l_c)$ with $l_c = 10$~m.

\subsubsection{Sensor Suite and State Estimation}

Each agent carries an IMU ($400$~Hz; gyroscope noise $5 \times 10^{-4}$~rad/s/$\sqrt{\text{Hz}}$; accelerometer noise $4 \times 10^{-3}$~m/s$^2$/$\sqrt{\text{Hz}}$), GPS ($10$~Hz; position noise $[0.02, 0.02, 0.05]$~m), barometer ($25$~Hz; $0.3$~m white noise, $0.2$~m correlated with $5$~s time constant), and a cable tension/direction sensor ($200$~Hz). State estimation uses a $15$-state error-state Kalman filter (ESKF) with multiplicative quaternion error, fusing IMU propagation at $400$~Hz with discrete GPS and barometer updates.

% ============================================================
\subsection{The N-Agnostic Constraint}
\label{sec:n_agnostic}
% ============================================================

Each agent~$i$ knows only:
\begin{itemize}[nosep]
  \item its own mass $m_Q$ and inertia $J$;
  \item its ESKF state estimate $(\hat{\bm{p}}_i, \hat{\bm{v}}_i, \hat{R}_i)$ and bias estimates;
  \item its cable tension $T_i$ and direction $\bm{q}_i$;
  \item a shared reference trajectory $\bm{p}_d(t)$ (pre-broadcast, evaluated locally).
\end{itemize}
It does \emph{not} know: the number of other agents, the payload mass, any other agent's state, or whether other agents are connected. This eliminates single points of failure and communication dependency, making the per-agent policy identical regardless of team size.

% ============================================================
\subsection{Baseline GPAC Architecture}
% ============================================================

The baseline GPAC architecture is a four-layer hierarchical controller per agent with separated timescales.

\subsubsection{Layer~1: Position and Anti-Swing Control (50~Hz)}

Computes the desired force:
\begin{equation}\label{eq:layer1_force}
  \bm{F}_{\mathrm{des},i} = \bm{F}_{\mathrm{fb}} + \bm{F}_{\mathrm{ff}} + \bm{F}_{\mathrm{cable}} + \bm{F}_{\mathrm{swing}} + \bm{F}_{\mathrm{eso}},
\end{equation}
where $\bm{F}_{\mathrm{fb}} = -K_p \bm{e}_p - K_d \bm{e}_v$ with gain matrices $K_p$, $K_d$; $\bm{F}_{\mathrm{ff}} = m_Q(\bm{a}_d + g\bm{e}_3)$; cable compensation $\bm{F}_{\mathrm{cable}} = \kappa_T(t) T_i \bm{q}_i$ with ramp $\kappa_T(t) = \min(1, (t - t_{\mathrm{taut}})/\Delta_{\mathrm{ramp}})$, $\Delta_{\mathrm{ramp}} = 2$~s; anti-swing $\bm{F}_{\mathrm{swing}} = k_q \bm{e}_q - k_\omega \bm{\omega}_q$ ($k_q = 4.0$~N/rad, $k_\omega = 2.0$~N$\cdot$s/rad); and ESO feedforward $\bm{F}_{\mathrm{eso}} = m_Q \hat{\bm{d}}_i$. The fault-tolerance evaluation (Sections~\ref{Section_IV:Fault-Tolerant_Control_Architecture}--\ref{Section_VI:Simulation_Results}) uses uniform scalar gains $k_p = 8.0$, $k_d = 8.0$, $k_i = 0.15$ (i.e., $K_p = k_p I$, $K_d = k_d I$) to enable the per-axis eigenvalue analysis in Section~\ref{sec:topology_invariance_theorem}.

\subsubsection{Layer~2: Geometric $\SOthree$ Attitude Control (200~Hz)}

Singularity-free tracking~\cite{lee2010geometric}:
\begin{equation}\label{eq:layer2_attitude}
  \bm{\tau}_i = -k_R \bm{e}_R - k_\Omega \bm{e}_\Omega + \bm{\Omega}_i \times (J\bm{\Omega}_i),
\end{equation}
with $\bm{e}_R = \tfrac{1}{2}\,\mathrm{vee}(R_d^\top R - R^\top R_d)$, $\bm{e}_\Omega = \bm{\Omega} - R^\top R_d \bm{\Omega}_d$, $k_R = 8.0$~N$\cdot$m, $k_\Omega = 1.5$~N$\cdot$m$\cdot$s. Almost-global exponential convergence holds on $\{\Psi_R < 2\} \subset \SOthree$.

\subsubsection{Layers~3--4 and Safety Filter (Not Active in This Study)}

The full GPAC baseline additionally includes a concurrent learning mass estimator~\cite{chowdhary2010concurrent} (Layer~3, 200~Hz), an extended state observer~\cite{han2009pid} (Layer~4, 500~Hz), and an ISSf-CBF safety filter~\cite{ames2017control,kolathaya2019input} (200~Hz, enforcing tension, cable-angle, tilt, and collision constraints). These are implemented in the Tether\_Lift codebase but are \emph{disabled} in the reported full-Drake runs (Remark~\ref{rem:system_under_test}) because the evaluation targets the minimal PID-plus-gravity-feedforward path. Their integration is discussed as future work in Section~\ref{Section_VII:Discussion}.

% ============================================================
\subsection{The Cable-Snap Fault Event}
\label{sec:cable_snap}
% ============================================================

At time $t_f$, cable~$j$ breaks instantaneously: $T_j(t) \to 0$ for $t \geq t_f$. Three coupled effects occur.

\emph{Payload acceleration:}
\begin{equation}\label{eq:payload_accel}
  \norm{\Delta \ddot{\bm{p}}_L} = \frac{T_j(t_f^-)}{m_L} \approx \frac{g}{N} \approx 3.27 \;\text{m/s}^2 \quad (N{=}3).
\end{equation}

\emph{Freed-agent acceleration:}
\begin{equation}\label{eq:freed_accel}
  \norm{\Delta \ddot{\bm{p}}_j} = \frac{T_j(t_f^-)}{m_Q} \approx \frac{g \cdot m_L}{N \cdot m_Q} \approx 6.54 \;\text{m/s}^2.
\end{equation}

\emph{Load-share jump:} from $m_L/N$ to $m_L/(N{-}1)$---a $50\%$ increase for $N{=}3$.

These effects propagate across two orders of magnitude in timescale:
\begin{center}
\small
\begin{tabular}{@{}lll@{}}
\toprule
Path & Mechanism & Timescale \\
\midrule
Rigid-body transmission & Force through payload & ${\sim}7.9$~ms \\
Elastic cable response & Spring-damper dynamics & $10$--$50$~ms \\
ESKF innovation & Velocity change detected & ${\sim}13$~ms \\
ESO convergence & Third-order settling & ${\sim}33$~ms \\
Load accel.\ filter & Low-pass differentiation & ${\sim}110$~ms \\
Cable wave propagation & $v_{\mathrm{wave}} \approx 7$~m/s & ${\sim}186$~ms \\
Adaptive divergence & Parameter estimate drift & ${\sim}300$~ms \\
\bottomrule
\end{tabular}
\end{center}

% ============================================================
\subsection{Two Fault Scenarios Under N-Agnosticism}
% ============================================================

\subsubsection{Scenario~A: Own Cable Breaks}

Agent~$i$ observes $T_i \to 0$ within one sample period ($\Delta t_d = 5$~ms). Detection is trivial (infinite SNR), and position deviation during detection is negligible:
\begin{equation}\label{eq:own_cable_deviation}
  \norm{\Delta \bm{p}_i(\Delta t_d)} \leq \frac{T_i^-}{2m_Q} \cdot \Delta t_d^2 \approx 0.082 \;\text{mm}.
\end{equation}
A threshold $T_i < T_{\min}$ (held for $0.12$~s debounce) triggers an escape maneuver: climbing $0.9$~m and displacing $1.8$~m horizontally within $1.5$~s, extending to $1.8$~m climb and $3.8$~m displacement within $3.0$~s. Direction is biased $0.18$ toward the formation tangent and $0.10$ toward current velocity, chosen to limit tilt to ${<}30$\textdegree\ during escape and ensure adequate vertical thrust margin.

\subsubsection{Scenario~B: Another Agent's Cable Breaks}

Agent~$i$ does not observe $T_j$ directly. The fault manifests only through payload-mediated perturbations. The tension-channel SNR is
\begin{equation}\label{eq:snr_ps}
  \mathrm{SNR}_T^{(B)} = \frac{\Delta T_i^{(\mathrm{imm})}}{\sigma_T} \approx \frac{4.2\;\text{N}}{3.0\;\text{N}} \approx 1.3,
\end{equation}
insufficient for reliable discrimination from wind gusts. Multi-channel analysis improves SNR to ${\sim}2$--$6$ but with $13$--$300$~ms delays. Scenario~B is the fundamentally difficult case and this paper's focus.

% ============================================================
\subsection{Why Conventional Fault Tolerance Cannot Apply}
\label{sec:why_fdir_fails}
% ============================================================

Under the N-agnostic constraint, conventional FDIR is \emph{architecturally incompatible} with Scenario~B: (i)~MMAE requires $N{+}1$ parallel filters, violating N-agnosticism; (ii)~tension-channel SNR~$\approx 1.3$ is indistinguishable from wind perturbations; (iii)~an N-agnostic agent has no mode parameter depending on~$N$ to reconfigure. These are structural consequences of the design choice. Any compatible fault-tolerance approach must work \emph{without} detection, isolation, or recovery.

% ============================================================
\subsection{Agent-Level Uncertainty Reformulation}
% ============================================================

From agent~$i$'s perspective, translational dynamics~\eqref{eq:quad_trans} can be rewritten as
\begin{equation}\label{eq:lumped_dynamics}
  \ddot{\bm{p}}_i = \frac{f_i}{m_Q} R_i \bm{e}_3 + g\bm{e}_3 + \bm{d}_i(t, \bm{x}),
\end{equation}
where $\bm{d}_i = (\bm{F}_i^{\mathrm{cable}} + \bm{F}_i^{\mathrm{wind}})/m_Q$ is the lumped acceleration uncertainty. This term encompasses cable tension (${\sim}6.5$~m/s$^2$), wind (${\sim}0.7$--$2.0$~m/s$^2$), and cable vibration (${\sim}0.3$~m/s$^2$). The controller never decomposes $\bm{d}_i$ into named components---this lumping enables topology invariance, since a cable snap changes the composition of $\bm{d}_i$ but not its mathematical role.

\begin{definition}[Nominal and Fault-Inclusive Disturbance Sets]\label{def:disturbance_sets_ps}
$\mathcal{W}_{\mathrm{nom}} = \{\bm{d} : \norm{\bm{d}}_\infty \leq \bar{d}_{\mathrm{nom}}\}$ bounds $\bm{d}_i$ under all-cables-intact operation. $\mathcal{W}_{\mathrm{FT}} \supset \mathcal{W}_{\mathrm{nom}}$ additionally bounds the worst-case cable-break transient. Both bounds are derived from system parameters. The dominant contribution is cable tension. Under quasi-static hover each cable carries $T_i \approx m_L g / N$, giving a cable-tension acceleration $d_{\mathrm{cable}} = m_L g / (N m_Q)$. Adding wind drag $d_{\mathrm{wind}} \leq C_D A_{\mathrm{ref}} \rho_{\mathrm{air}} v_{\mathrm{wind,max}}^2 / (2 m_Q) \approx 0.7$~m/s$^2$ and cable vibration $d_{\mathrm{vib}} \leq k_s \varepsilon_{\max} / (m_Q (n_b+1)) \approx 0.3$~m/s$^2$ yields
\begin{equation}\label{eq:sigma_nom_derivation}
  \bar{d}_{\mathrm{nom}} = \frac{m_L g}{N m_Q} + d_{\mathrm{wind}} + d_{\mathrm{vib}} \approx \frac{3.0 \times 9.81}{3 \times 1.5} + 1.0 \approx 7.5 \;\text{m/s}^2.
\end{equation}
Post-fault, the surviving $N{-}k$ agents share the load, so the cable-tension contribution increases to $m_L g / ((N{-}k) m_Q)$:
\begin{equation}\label{eq:sigma_ft_derivation}
  \bar{d}_{\mathrm{FT}} = \frac{m_L g}{(N{-}k) m_Q} + d_{\mathrm{wind}} + d_{\mathrm{vib}}.
\end{equation}
For $N{=}3$, $k{=}1$: $\bar{d}_{\mathrm{FT}} \approx 9.81 + 1.0 \approx 10.8$~m/s$^2$. With safety margin for elastic transients, we use the rounded values $\bar{d}_{\mathrm{nom}} \approx 8$~m/s$^2$ and $\bar{d}_{\mathrm{FT}} \approx 12$~m/s$^2$. In the cable-dominated regime ($d_{\mathrm{cable}} \gg d_{\mathrm{wind}} + d_{\mathrm{vib}}$), the cost ratio simplifies to a closed-form expression:
\begin{equation}\label{eq:rho_closed_form}
  \rho_{\mathrm{FT}} \approx \frac{\bar{d}_{\mathrm{FT}}}{\bar{d}_{\mathrm{nom}}} \approx \frac{N}{N{-}k},
\end{equation}
which gives $\rho_{\mathrm{FT}} = 3/2 = 1.5$ for $N{=}3$, $k{=}1$; $5/3 = 1.67$ for $N{=}5$, $k{=}2$; and $7/4 = 1.75$ for $N{=}7$, $k{=}3$. These are conservative upper bounds; the empirical values ($1.26$--$1.39$) are lower because the fault disturbance is transient rather than sustained (Section~\ref{sec:topology_invariance_theorem}).
\end{definition}

% ============================================================
\subsection{Formal Problem Statement}
% ============================================================

Given the system on manifold~\eqref{eq:config_manifold}, the N-agnostic constraint (Section~\ref{sec:n_agnostic}), the cable-snap fault model (Section~\ref{sec:cable_snap}), the lumped uncertainty~\eqref{eq:lumped_dynamics} with disturbance sets (Definition~\ref{def:disturbance_sets_ps}), and the architectural incompatibility of detection-based FDIR (Section~\ref{sec:why_fdir_fails}), the paper addresses two problems.

\begin{enumerate}[label=\textbf{P\arabic*.},leftmargin=*]
  \item \textbf{Topology-Invariant Tracking.} Identify sufficient conditions under which the per-agent tracking error satisfies
  \begin{equation}\label{eq:problem_bound}
    \norm{\bm{e}_i(t)}_\infty \leq \gamma_1(T_s, \omega, \bar{d}_{\mathrm{FT}}) \quad \forall\, t \geq 0
  \end{equation}
  for \emph{all cable topologies}---where $\gamma_1$ depends on controller bandwidth and disturbance size, not on cable count or identity---provided $\bm{d}_i \in \mathcal{W}_{\mathrm{FT}}$ and thrust authority is adequate.

  \item \textbf{Fault-Tolerance Cost.} Derive the tradeoff
  \begin{equation}\label{eq:problem_rho}
    \rho_{\mathrm{FT}} = \frac{\gamma_1(T_s, \omega, \bar{d}_{\mathrm{FT}})}{\gamma_1(T_s, \omega, \bar{d}_{\mathrm{nom}})} > 1
  \end{equation}
  between nominal tracking and implicit fault tolerance, providing a design criterion for choosing between detection-free tolerance and detection-based recovery.
\end{enumerate}

\textbf{P1} is addressed by a structural analysis (Proposition~\ref{thm:topology_invariance}) and comprehensive full-Drake simulation evidence. \textbf{P2} is addressed by both steady-state and transient-aware analysis, validated empirically.
